% !TeX root = ../../Book2.tex
\chapter*{后记}
\addcontentsline{toc}{chapter}{后记}

回望本卷，许多证明都在做同一件具体的事：先选择一个能够保存的信息，再让其余部分趋于极限。Fourier 截断保存频率，能量估计保存范数，紧性保存收敛子列，弱形式则保存分部积分关系。选择不同，得到的解类与收敛意义也不同；它们不能只凭同一个形式公式彼此替换。

\ProseParagraphBreak

表示公式、唯一性与正则性回答不同的问题。一条积分公式可以构造解，能量估计说明同一数据是否还允许其他解，正则性则使积分意义的方程重新具有经典导数。算子的定义域又把这些性质同边界条件联系起来。沿着这些联系，经典解、弱解、分布解和半群解成为观察同一个方程的不同方式。

\ProseParagraphBreak

高阶调和分析把这种联系推进到尺度与几何。局部核估计需要在许多频带上求和，曲率带来的衰减需要与波包的集中相比较，相位又把奇性的位置与传播方向连接起来。继续研究限制问题或微局部分析时，可以沿正文所列文献追踪这些构造，观察每一层尺度怎样传递已有的估计。

\ProseParagraphBreak

合上本卷之前，可以重做三组问题。先从高斯函数推导 Fourier 变换、热核与逼近恒等，检查三种推导中的常数是否一致；再把一个 Dirichlet 问题分别写成能量极小化、双线性型和弱方程，指出存在、唯一与正则性各用了什么；最后选择一个乘子或振荡核，先作缩放分析，再判断绝对可积、抵消与曲率分别提供哪一种估计。这些问题覆盖的页数不多，却能检验是否真正连接起了本卷的方法。

\ProseParagraphBreak

后续道路会逐渐分开。非线性方程需要更多紧性、稳定性和解的选择原则；调和分析需要更精细的局部化、几何分解与端点技巧；谱论与微局部分析还要认真处理无界算子和相空间传播。特殊函数的显式计算在这些方向中仍然有用，但它们通常只是模型，不能代替一般估计。

\ProseParagraphBreak

一本教材留下的不应只有一串可引用的结果。更希望留下的是一些可以自己检查的动作：写出测试函数，算一次尺度，追踪一个边界项，尝试删去一条假设，再找出失败的位置。当一个新问题没有现成公式时，这些动作往往就是开始证明的地方。

\bigskip
\begin{flushright}
R. EverStarine

2026年10月11日
\end{flushright}
